% 09_body.tex — Day2 산출물 4/5 (⑨) 원가 기준선과 예비비 구조 (빈 템플릿 / 완성 예시 공용 본문)
% 드라이버: 09_원가기준선_빈템플릿.tex (\wbblanktrue) / 09_원가기준선_완성예시.tex (\wbblankfalse)
% 내용 출처: PM트랙/Day2_워크북.md §4.3(항목 가이드) §4.4(빈 템플릿) §4.5(온담 P1 완성 예시본)
\documentclass[10pt,a4paper]{article}
\usepackage[a4paper,margin=16mm,top=14mm,bottom=20mm]{geometry}
\def\DocNum{4}
\def\DocTitle{원가 기준선과 예비비 구조}
\def\DocSub{⑨ Cost Baseline \& Reserves}
\def\DocVariant{\B{빈 템플릿}{온담 P1 완성 예시}}
\usepackage{wbstyle}
\begin{document}

\docheader
 {\Bg{[정식 명칭과 코드]}{차세대 주문·재고 통합 플랫폼 구축 (ONE ONDAM P1)}}
 {\Bg{[v0.x / 승인 게이트]}{v0.9 / 2026-03-27, G1 2026-04-30 승인 예정}}
 {\Bg{[PM 실명 / PMO 기준선 담당]}{P1 PM / P1 PMO(기준선)\rd{7mm}}}
 {\Bg{[스폰서 실명]}{정미란 CFO (검토: 재경팀장, 이재현)\rd{7mm}}}

% ------------------------------------------------------------------ 1
\secbar{1. 문서 정보와 견적 등급}
\guide{기준선 버전, 승인 게이트, 통화·단위(억 원, 소수 자릿수), 견적 등급과 근거(어느 설계 성숙도에서 만든 숫자인가), 상위 문서. 고정가 부분과 추정 부분을 구분해 어느 부분이 등급의 대상인지 밝힌다. 등급은 AACE Class 표기를 빌리되 “정의도 n\%, 정확도 범위는 별도 산정”으로 — \textbf{Class에서 정확도(“Class 3이므로 −20/+30\%”)를 자동으로 끌어오지 않는다.} 등급 없이 숫자만 적으면 소수 첫째 자리가 정확도인 것처럼 읽힌다.}
{\footnotesize
\begin{tabularx}{\linewidth}{|L{26mm}|Y|L{24mm}|Y|}\hline
\hd{항목} & \hd{내용} & \hd{항목} & \hd{내용} \\ \hline
\B{\rd{9mm}}{}버전 / 승인 & \Bg{[v0.x / 게이트, 날짜, 승인 주체]}{v0.9 / G1 2026-04-30 (RACI 14행, 이사회 A)} & 단위 & \Bg{[억 원, 소수 n자리]}{억 원, 소수 둘째 자리} \\ \hline
\B{\rd{16mm}}{}고정가 부분 & \Bg{[금액, 계약 근거 — 등급 대상 아님]}{SI 68.20 (FP 12,400 × 55만 원, P1 계약서 부속서 1) — 등급 대상 아님} & 추정 부분·등급 & \Bg{[금액 — Class n 상당, 정의도, 정확도 범위 산정 방식]}{77.60 — 설계 스프린트 S6 종료 기준 “AACE Class 3 상당” [추가 설정]. 정의도: 아키텍처 확정, 인프라 사이징 실측 전(클라우드 3.0 유보 사유). 정확도 범위: 온담 과거 IT 투자 실적 분포 부재로 미산정, 컨틴전시 배정 논리(항목 6)로 대체} \\ \hline
\B{\rd{9mm}}{}상위 문서 & \Bg{[BC §5·§6, 헌장 §8, WBS 항목 9, 리스크 등록부]}{BC 요약 v0.9 §5·§6, 헌장 v1.0 §8, WBS v0.9 항목 9, 리스크 등록부 v1.0(⑩)} & 작성 / 검토 / 승인 & \Bg{[실명 / 실명 / 실명]}{P1 PM, P1 PMO(기준선) / 재경팀장, 이재현 / 정미란} \\ \hline
\end{tabularx}}

% ------------------------------------------------------------------ 2
\secbar[90mm]{2. 3층 예산 표}
\guide{승인 예산 / 집행 한도 / 계획 원가(기준선)의 금액·정의·통제 주체·결재선과 각 층의 \textbf{변경 절차}. 검산식: 계획 원가 + 컨틴전시 = 집행 한도, 집행 한도 + 관리예비비 = 승인 예산. 컨틴전시는 기준선 위에 있고 집행될 때 기준선으로 이동한다 — PV 곡선에 섞지 말 것.}
{\footnotesize
\begin{tabularx}{\linewidth}{|L{28mm}|C{16mm}|Y|L{36mm}|L{36mm}|}\hline
\hd{층} & \hd{금액} & \hd{정의} & \hd{통제 주체 / 결재선} & \hd{변경 절차} \\ \hline
\ifwbblank
\rd{10mm}승인 예산 & & \g{[의결 근거]} & \g{[이사회]} & \\ \hline
\rd{10mm}관리예비비 & & \g{[용도: 미식별·프로그램 계층 원인·유보 복귀]} & \g{[스폰서 A, 재경 C]} & \g{[항목 7 절차]} \\ \hline
\rd{10mm}집행 한도 & & \g{= 승인 − 관리예비비} & & \\ \hline
\rd{10mm}컨틴전시 & & \g{[배정 \hspace{6mm} + 미배정 풀 \hspace{6mm}]} & \g{[PM A, 재경 C, CCB I]} & \g{[단일 건 기준 금액, 사내 지침]} \\ \hline
\rd{10mm}\textbf{계획 원가 (기준선)} & & \g{= 집행 한도 − 컨틴전시 = WBS L2 합 = PV 누적} & & \\ \hline
\else
승인 예산 & 168.00 & 이사회 의결 계약·발주 한도 (P-08) & 프로그램 이사회 & 프로그램 재정의 \\ \hline
관리예비비 & 12.00 & 프로그램 24억 중 P1 유보분. 미식별 리스크, 프로그램·스폰서 계층 원인 리스크(R-01·R-04), 유보 조정 항목 복귀 & 정미란 (RACI 20행 A), 재경팀장 C & 항목 7 절차. 집행 시 기준선 갱신 \\ \hline
집행 한도 & 156.00 & 승인 예산 − 관리예비비. 스폰서 결재 없이 쓸 수 있는 상한 & P1 PM & — \\ \hline
컨틴전시 & 10.20 & 알려진 리스크 잔여 EMV 배정 4.60 + 미배정 풀 5.60. 집행 시 기준선으로 이동(기준선 변경 아님) & P1 PM (RACI 19행 A), 재경팀장 C(한도 확인), CCB I(월 보고), 윤태호 사본 & 5억 초과 단일 건 CFO 병행, 5,000만 원 이상 내부통제 지침 \\ \hline
\textbf{계획 원가 (원가 기준선)} & \textbf{145.80} & 집행 한도 − 컨틴전시 = WBS L2 합 = PV 누적 & P1 PM & 범위 변경(CCB) 또는 관리예비비 집행 시 갱신 \\ \hline
\fi
\end{tabularx}}
\B{\small\g{검산: \hspace{12mm} + \hspace{12mm} = 집행 한도 \hspace{12mm}, \hspace{12mm} + \hspace{12mm} = 승인 예산 \hspace{12mm}}\par}{\small 검산: 145.80 + 10.20 = 156.00, 156.00 + 12.00 = 168.00.\par}

% ------------------------------------------------------------------ 3 (가로)
\landscapeon
\secbar{3. 승인 내역 → 기준선 조정표}
\guide{승인 내역 항목과 기준선 항목의 차이를 항목별로 어떻게 조정했는지, 조정의 성격(타 프로젝트 이관 / 유보 / 범위 결정 / 기간 정합), 근거·승인자, 조정이 깨질 때(되돌아올) 조건. \textbf{조정 합계 = 관리예비비}여야 한다. 조정은 “숫자를 줄이는 것”이 아니라 “누가 언제 내는 돈인지를 다시 정하는 것”이다. 고정가 계약 부분을 줄이려 하지 말 것.}
{\footnotesize
\begin{longtable}{|L{36mm}|C{16mm}|C{14mm}|C{16mm}|L{70mm}|L{46mm}|L{50mm}|}\hline
\hd{항목} & \hd{승인 내역} & \hd{조정} & \hd{기준선} & \hd{조정 성격 (이관 / 유보 / 범위 결정 / 기간 정합) — 내용} & \hd{근거·승인자} & \hd{되돌아올 조건} \\ \hline \endhead
\ifwbblank
\rd{11mm}\g{[항목 1 — 고정가]} & & \g{0} & & \g{고정가 — 조정 불가} & \g{[계약]} & \g{—} \\ \hline
\rd{11mm} & & & & & & \\ \hline
\rd{11mm} & & & & & & \\ \hline
\rd{11mm} & & & & & & \\ \hline
\rd{11mm} & & & & & & \\ \hline
\rd{11mm} & & & & & & \\ \hline
\rd{11mm} & & & & & & \\ \hline
\rd{11mm} & & & & & & \\ \hline
\rd{9mm}\textbf{합계 (컨틴전시 제외)} & & \g{[= −관리예비비]} & \g{[= 계획 원가]} & & & \\ \hline
\else
SI 개발 용역 & 68.20 & 0 & 68.20 & 고정가 — 조정 불가 & 계약 & — \\ \hline
상용SW 라이선스·유지보수 & 23.40 & −4.00 & 19.40 & \textbf{P5 이관} — 3년 유지보수 6.0 중 G4(2027-05) 이후 발생하는 2·3년차 4.0(연 2.0)은 P5 운영비 27억의 항목. 라이선스 17.4 + 1년차 유지보수 2.0 & 박준영·재경팀장 합의 2026-03, 정미란 승인 & P5 예산에서 거부되면 관리예비비 요청 \\ \hline
인프라 & 23.90 & −3.00 & 20.90 & \textbf{유보} — 클라우드 3년 선약정 14.6을 설계 스프린트 실측 사이징 후 2년 선약정 + 1년 종량으로 재구성, 기준선 11.6. 차액 3.0은 사이징 초과 시 스폰서 결재로 복귀 & 박세연, 강현도 / 정미란 & 부하시험(2027-01) 결과 피크 5,100 케이스/h 기준 초과 시 \\ \hline
POS 단말 & 8.60 & 0 & 8.60 & — & & \\ \hline
데이터 전환 용역 & 11.40 & −2.00 & 9.40 & \textbf{범위 결정} — 5년 초과 이력(약 2.6TB) 아카이브만, 신 플랫폼 이관 제외 (범위 기술서 §5 [신규] 가정) & 박세연, 재경팀장(보존 요건 확인) — 헌장 v1.1 & 세무·감사 보존 요건이 이관을 요구하면 복귀 \\ \hline
통합·성능시험·감리 & 9.70 & 0 & 9.70 & — & & \\ \hline
발주사 PMO·형상·품질 & 12.60 & −3.00 & 9.60 & \textbf{기간 정합} — 3명 × 24개월(10.8) + 도구 1.8 중 P1 기간 17개월분 7.65 + 도구 1.8 = 9.45 → 9.60(반올림·인수인계 0.15). 잔여 7개월 3.0은 G4 이후 하자보수기 PMO로 유보 & 프로그램 PMO장, 정미란 & G4 이후 하자보수기 착수 시 \\ \hline
\textbf{합계 (컨틴전시 제외)} & \textbf{157.80} & \textbf{−12.00} & \textbf{145.80} & 조정 합계 = 관리예비비 12.0 [조정 내역은 추가 설정] & & \\ \hline
\fi
\end{longtable}}

% ------------------------------------------------------------------ 4 (가로)
\secbar[70mm]{4. WBS 레벨 2 배분}
\guide{⑥ WBS 항목 9와 같은 교차표. 원가 기준선 관점에서 “어느 L2가 어느 원가 항목을 얼마나 갖는가”를 확인하고 합계 = 계획 원가를 검산한다. \textbf{두 문서의 표가 다르면 어느 쪽이 틀렸는지 찾는다.} 배분 없이 항목 합계만 있는 기준선은 변경의 원가 영향을 붙일 노드가 없다.}
{\footnotesize
\begin{tabularx}{\linewidth}{|Y|C{18mm}|C{24mm}|C{18mm}|C{18mm}|C{18mm}|C{20mm}|C{20mm}|C{18mm}|C{22mm}|}\hline
\hd{L2} & \hd{FP} & \hd{SI (FP × 단가)} & \hd{상용SW} & \hd{인프라} & \hd{단말} & \hd{데이터 전환} & \hd{시험·감리} & \hd{PMO} & \hd{합계 (억)} \\ \hline
\ifwbblank
\rd{8mm}1.1 & & & & & & & & & \\ \hline
\rd{8mm}1.2 & & & & & & & & & \\ \hline
\rd{8mm}1.3 & & & & & & & & & \\ \hline
\rd{8mm}1.4 & & & & & & & & & \\ \hline
\rd{8mm}1.5 & & & & & & & & & \\ \hline
\rd{8mm}1.6 \g{[프로젝트 관리]} & \g{—} & \g{—} & & & & & & & \\ \hline
\rd{8mm}\textbf{합계} & \g{[= 총 FP]} & & & & & & & & \g{[= 계획 원가]} \\ \hline
\else
1.1 POS·매장 운영 & 3,100 & 17.05 & — & — & 8.60 & — & — & — & \textbf{25.65} \\ \hline
1.2 주문 허브 & 2,900 & 15.95 & — & — & — & — & — & — & \textbf{15.95} \\ \hline
1.3 통합 재고 & 2,300 & 12.65 & — & — & — & — & — & — & \textbf{12.65} \\ \hline
1.4 가맹 발주 포털 & 1,500 & 8.25 & — & — & — & — & — & — & \textbf{8.25} \\ \hline
1.5 통합 계층·전환·공통 기반 & 2,600 & 14.30 & 19.40 & 20.90 & — & 9.40 & — & — & \textbf{64.00} \\ \hline
1.6 프로젝트 관리·통합 검증 & — & — & — & — & — & — & 9.70 & 9.60 & \textbf{19.30} \\ \hline
\textbf{합계} & \textbf{12,400} & \textbf{68.20} & \textbf{19.40} & \textbf{20.90} & \textbf{8.60} & \textbf{9.40} & \textbf{9.70} & \textbf{9.60} & \textbf{145.80} \\ \hline
\fi
\end{tabularx}}
\B{}{\small ⑥ WBS 항목 9와 동일. 합계 145.80 = 계획 원가 (항목 2).\par}

% ------------------------------------------------------------------ 5 (가로)
\clearpage
\secbar{5. 월별 PV 곡선 (성과측정기준선, PMB)}
\guide{프로젝트 기간의 월별 계획 가치(PV)와 누적, 원가 항목별 시간 배분 규칙(SI = 스프린트 가중, 라이선스 = 인도 시점, PMO = 균등 …), 누적 곡선의 형태와 이유. \textbf{PV는 지급 시점이 아니라 작업이 예정된 시점의 가치}다 — 검수 후 지급 일정을 PV로 쓰면 곡선이 뒤로 몰려 EVM이 무의미하다. 균등 배분(총액 ÷ 개월)으로 S-curve를 만들지 말고, 컨틴전시를 섞지 말 것. 연차 집행 계획(BC §5)과 대조해 정합을 확인한다.}
\noindent\textbf{\small 배분 규칙} \g{(항목별 한 줄)}\par\vspace{1mm}
\begin{fieldbox}
\ifwbblank
\g{SI: \hrulefill}\par\vspace{2.5mm}\g{상용SW: \hrulefill}\par\vspace{2.5mm}\g{인프라: \hrulefill}\par\vspace{2.5mm}\g{단말: \hrulefill}\par\vspace{2.5mm}\g{데이터 전환: \hrulefill}\par\vspace{2.5mm}\g{시험·감리: \hrulefill}\par\vspace{2.5mm}\g{PMO: \hrulefill}\par
\else
\small [추가 설정] \textbf{SI} = 스프린트 시작 월에 배정, 가중치 설계 0.7 / 구현 1.1 / 통합·시험 1.2 (6×0.7 + 20×1.1 + 10×1.2 = 38.2단위, 단위당 1.785) → 설계 스프린트 1.25, 구현 1.96, 통합 2.14. \textbf{상용SW} = 라이선스 인도 시점(3월 8.0 API GW·MDM 1차, 4월 6.0 MDM 2차·APM, 5월 3.4 SAP 애드온) + 1년차 유지보수 월 0.2 × 10(2026-06~2027-03). \textbf{인프라} = 선약정 6.0(4월) + 클라우드 월 0.4 × 14(2026-04~2027-05) + 망분리·보안 3.0(6월)·2.1(7월) + DR 2.1(12월)·2.1(1월). \textbf{단말} = 입고 3.0(12월)·4.06(1월) + 설치 0.8(1월)·0.74(2월). \textbf{데이터 전환} = 매핑~리허설 10개월 0.6→1.2→0.8. \textbf{시험·감리} = 감리 G1 0.8·G2 1.4·G3 1.4·G4 0.6, 보안진단 1.0(8월)·1.4(1월), 부하시험 1.5(1월)·1.6(2월). \textbf{PMO} = 월 0.56 × 16 + 0.64.
\fi
\end{fieldbox}
\B{\clearpage\noindent\textbf{\small 5. 월별 PV 표 (계속)} \g{— 월 = 프로젝트 기간 전체, 마지막 행 합계 = 계획 원가}\par\vspace{1mm}}{}
{\footnotesize
\begin{longtable}{|L{22mm}|C{18mm}|C{18mm}|C{18mm}|C{18mm}|C{18mm}|C{20mm}|C{18mm}|C{22mm}|C{24mm}|C{20mm}|}\hline
\hd{월} & \hd{SI} & \hd{상용SW} & \hd{인프라} & \hd{단말} & \hd{전환} & \hd{시험·감리} & \hd{PMO} & \hd{월 PV} & \hd{누적 PV} & \hd{누적 \%} \\ \hline \endhead
\ifwbblank
\rd{4.5mm}\g{[YYYY-MM]} & & & & & & & & & & \\ \hline
\rd{4.5mm} & & & & & & & & & & \\ \hline
\rd{4.5mm} & & & & & & & & & & \\ \hline
\rd{4.5mm} & & & & & & & & & & \\ \hline
\rd{4.5mm} & & & & & & & & & & \\ \hline
\rd{4.5mm} & & & & & & & & & & \\ \hline
\rd{4.5mm} & & & & & & & & & & \\ \hline
\rd{4.5mm} & & & & & & & & & & \\ \hline
\rd{4.5mm} & & & & & & & & & & \\ \hline
\rd{4.5mm} & & & & & & & & & & \\ \hline
\rd{4.5mm} & & & & & & & & & & \\ \hline
\rd{4.5mm} & & & & & & & & \g{[연차 누적 = BC 연차 계획?]} & & \\ \hline
\rd{4.5mm} & & & & & & & & & & \\ \hline
\rd{4.5mm} & & & & & & & & & & \\ \hline
\rd{4.5mm} & & & & & & & & & & \\ \hline
\rd{4.5mm} & & & & & & & & & & \\ \hline
\rd{4.5mm} & & & & & & & & & \g{[= 계획 원가]} & 100\% \\ \hline
\rd{4.5mm}\textbf{합계} & & & & & & & & \g{[= 계획 원가]} & & \\ \hline
\else
2026-01 & 2.50 & — & — & — & — & — & 0.56 & 3.06 & 3.06 & 2.1\% \\ \hline
2026-02 & 2.50 & — & — & — & — & — & 0.56 & 3.06 & 6.12 & 4.2\% \\ \hline
2026-03 & 4.46 & 8.00 & — & — & — & — & 0.56 & 13.02 & 19.14 & 13.1\% \\ \hline
2026-04 & 3.93 & 6.00 & 6.40 & — & — & 0.80 & 0.56 & 17.69 & 36.83 & 25.3\% \\ \hline
2026-05 & 3.93 & 3.40 & 0.40 & — & 0.60 & — & 0.56 & 8.89 & 45.72 & 31.4\% \\ \hline
2026-06 & 3.93 & 0.20 & 3.40 & — & 0.80 & — & 0.56 & 8.89 & 54.61 & 37.5\% \\ \hline
2026-07 & 3.93 & 0.20 & 2.50 & — & 0.90 & — & 0.56 & 8.09 & 62.69 & 43.0\% \\ \hline
2026-08 & 5.89 & 0.20 & 0.40 & — & 1.00 & 1.00 & 0.56 & 9.05 & 71.74 & 49.2\% \\ \hline
2026-09 & 3.93 & 0.20 & 0.40 & — & 1.00 & — & 0.56 & 6.09 & 77.83 & 53.4\% \\ \hline
2026-10 & 3.93 & 0.20 & 0.40 & — & 1.20 & — & 0.56 & 6.29 & 84.12 & 57.7\% \\ \hline
2026-11 & 3.93 & 0.20 & 0.40 & — & 1.20 & 1.40 & 0.56 & 7.69 & 91.81 & 63.0\% \\ \hline
2026-12 & 3.93 & 0.20 & 2.50 & 3.00 & 1.00 & — & 0.56 & 11.19 & \textbf{103.00} & 70.6\% \\ \hline
2027-01 & 4.28 & 0.20 & 2.50 & 4.86 & 0.90 & 2.90 & 0.56 & 16.20 & 119.20 & 81.8\% \\ \hline
2027-02 & 4.28 & 0.20 & 0.40 & 0.74 & 0.80 & 3.00 & 0.56 & 9.98 & 129.19 & 88.6\% \\ \hline
2027-03 & 6.43 & 0.20 & 0.40 & — & — & — & 0.56 & 7.59 & 136.77 & 93.8\% \\ \hline
2027-04 & 4.28 & — & 0.40 & — & — & — & 0.56 & 5.24 & 142.02 & 97.4\% \\ \hline
2027-05 & 2.14 & — & 0.40 & — & — & 0.60 & 0.64 & 3.78 & \textbf{145.80} & 100.0\% \\ \hline
\textbf{합계} & \textbf{68.20} & \textbf{19.40} & \textbf{20.90} & \textbf{8.60} & \textbf{9.40} & \textbf{9.70} & \textbf{9.60} & \textbf{145.80} & & \\ \hline
\fi
\end{longtable}}
\B{\small\g{곡선의 형태와 이유 (봉우리 = 조달·단말 인도 시점): \hrulefill}\par\vspace{2mm}\g{연차 집행 계획(BC §5)과의 대조: \hrulefill}\par}{\small 월 값은 소수 둘째 자리 반올림, 누적은 원값 기준이라 손으로 더하면 ±0.01 차이가 날 수 있다. 형태: 2026-03~04 라이선스·선약정 봉우리(누적 25\%까지 4개월), 이후 월 6~9억의 완만한 구현 구간, 2026-12~2027-02 단말·DR·시험 봉우리(누적 70→89\%), 2027-03~05 꼬리(안정화·이관). 2026년 누적 103.00 vs BC 연차 집행 계획 102.0 [추가 설정] — 차이 1.0은 PV(작업 예정)와 현금 집행(검수 후 지급)의 시점 차이, 정합. \textbf{CFO의 1차연도 148억 요구(R-04)가 확정되면 3~4월 봉우리(라이선스 14.0·선약정 6.0)가 첫 번째 타격 지점이다.}\par}

% ------------------------------------------------------------------ 6 (가로)
\clearpage
\secbar{6. 컨틴전시 배정표}
\guide{컨틴전시를 리스크 ID별 잔여 EMV(⑩)에 따라 배정한 표와 미배정 풀의 용도·규모·재배정 주기. 배정 합 + 미배정 풀 = 컨틴전시. 미배정 풀의 근거는 프로젝트의 관리 단위(스프린트 1회 비용)로 설명한다. \textbf{관리예비비 대상 리스크(원인·결정이 스폰서·프로그램 계층)는 여기 포함하지 않고 별도 표시} — 컨틴전시로 막으려 하면 그 리스크가 실현될 때 PM은 다른 모든 리스크에 무방비가 된다. “총액 × n\%”로 끝내지 말 것.}
{\footnotesize
\begin{longtable}{|L{74mm}|C{20mm}|C{16mm}|L{62mm}|L{44mm}|L{32mm}|}\hline
\hd{리스크 ID — 사건} & \hd{잔여 EMV (⑩)} & \hd{배정액} & \hd{트리거 (집행 개시)} & \hd{집행 결재} & \hd{자금원} \\ \hline \endhead
\ifwbblank
\rd{7mm}\g{R-[\ ] [사건 요약]} & & & \g{[관측 가능한 신호·시점]} & \g{[PM / 기준 금액 초과 시 병행 결재]} & \g{컨틴전시} \\ \hline
\rd{7mm} & & & & & \\ \hline
\rd{7mm} & & & & & \\ \hline
\rd{7mm} & & & & & \\ \hline
\rd{7mm} & & & & & \\ \hline
\rd{7mm} & & & & & \\ \hline
\rd{7mm} & & & & & \\ \hline
\rd{7mm} & & & & & \\ \hline
\rd{7mm} & & & & & \\ \hline
\rd{7mm} & & & & & \\ \hline
\rd{7mm} & & & & & \\ \hline
\rd{7mm}\textbf{배정 소계} & & & & & \\ \hline
\rd{7mm}\textbf{미배정 풀} & — & & \g{[용도 · 규모 근거(스프린트 n회분) · 재배정 주기]} & & \\ \hline
\rd{7mm}\textbf{합계} & & \g{[= 컨틴전시]} & & & \\ \hline
\rd{7mm}\textit{관리예비비 대상 (별도)} & \g{[ID: EMV, …]} & — & \g{[원인·결정이 프로그램·스폰서 계층]} & \g{[스폰서]} & \g{관리예비비} \\ \hline
\else
R-02 가맹 부속서 지연 → 분리 배포·재작업 & 0.60 & 0.60 & 2026-07 동의율 중간 집계 60\% 미만 & P1 PM, 재경 확인 & 컨틴전시 \\ \hline
R-03 3PL 거부 → 폴백 확정·재시험 & 0.48 & 0.50 & 2026-06-30 미타결 & P1 PM & 컨틴전시 \\ \hline
R-05 겸임 인력 → 대기 비용 청구 & 0.60 & 0.60 & 변경협의체 대기 비용 판정 & P1 PM (5,000만 원 이상 건 CFO) & 컨틴전시 \\ \hline
R-06 OD-POS 개발자 이탈 → 외주 분석 & 0.30 & 0.30 & 퇴직 의사 표명 또는 문서화 진도 2026-03 50\% 미만 & P1 PM & 컨틴전시 \\ \hline
R-07 요구 집중 → FP 재산정 & 0.40 & 0.40 & 월 순변동 ±35건 초과 & P1 PM & 컨틴전시 \\ \hline
R-08 앱 벤더 계약 외 → P1측 연동 재작업 & 0.60 & 0.60 & 벤더 API 규격 회신 4주 지연 & P1 PM (벤더 견적 자체는 앱 유지보수 계약, P1 예산 외) & 컨틴전시 \\ \hline
R-09 별도 법인 규제 주체 → 공장 연동 재설계 & 0.30 & 0.30 & P4 위탁 구조 2026-05 미확정 & P1 PM & 컨틴전시 \\ \hline
R-10 G2 후 마스터 변경 → 재매핑 & 0.38 & 0.40 & 정본 판정 후 마스터 변경 CCB 승인 & P1 PM → 3.8억 전액 시 CFO 병행 & 컨틴전시 (초과분 관리예비비) \\ \hline
R-11 SAC 미합의 → 하자보수 연장 & 0.40 & 0.40 & G1 SAC 초안 미합의 & P1 PM & 컨틴전시 \\ \hline
R-12 편익 분모·실사 비협조 → 재실사 & 0.20 & 0.20 & M+3 분모 미확정, 실사 협조 거부 & P1 PM & 컨틴전시 \\ \hline
R-13 구형 단말 → 롤백·재방문 & 0.32 & 0.30 & 단말 조사 비호환 10\% 초과 (2026-04-30) & P1 PM & 컨틴전시 \\ \hline
\textbf{배정 소계} & \textbf{4.58} & \textbf{4.60} & & & \\ \hline
\textbf{미배정 풀} & — & \textbf{5.60} & 월 동결 시 신규 식별 리스크, 스프린트 재계획·재작업 — 스프린트 4회분(1.41 × 4 = 5.64)에 해당 [추가 설정]. 분기 재배정 & P1 PM, CCB 월 보고 & 컨틴전시 \\ \hline
\textbf{합계} & & \textbf{10.20} & & & \\ \hline
\textit{관리예비비 대상 (별도)} & R-01 1.20 + R-04 1.00 = \textbf{2.20} & — & 원인·결정이 프로그램·스폰서 계층. T0 리드타임 확정 / 자금계획 확정 & 정미란 & 관리예비비 12.0 \\ \hline
\fi
\end{longtable}}
\landscapeoff

% ------------------------------------------------------------------ 7
\secbar[90mm]{7. 관리예비비 집행 절차}
\guide{관리예비비의 용도 정의(미식별 리스크, 프로그램 계층 원인 리스크, 유보 조정 항목의 복귀), 집행 요청 문서, 결재선, 시한, 기록·보고를 5~7단계로. 헌장 §12의 예외 보고 형식(원인·영향·대안 3·권고)을 집행 요청 형식으로 재사용한다. \textbf{관리예비비가 집행되면 그만큼 기준선이 늘어난다}(재기준선 아님, 승인된 변경) — 갱신하지 않으면 EVM이 어긋난다.}
\begin{fieldbox}
\ifwbblank
\foreach \i/\t in {1/요청 — 누가·어떤 형식·대상,2/검토 — 누가·무엇을·시한,3/결재 — 누가·시한·기준 금액 초과 시 병행,4/보고 — 회의체·사본,5/기준선 갱신 — BL 번호·PV 곡선·집행 한도,6/기록·사본 — 어디에}{\g{\i. [\t]} \rule{0pt}{9mm}\hrulefill\par}
\else
\begin{enumerate}
\item \textbf{요청} — P1 PM이 예외 보고 형식(원인·영향·대안 3·권고, 헌장 §12)으로 집행 요청서 작성. 대상: 관리예비비 대상 리스크 실현(R-01·R-04), 미식별 리스크, 유보 조정 항목 복귀(클라우드 3.0, PMO 3.0, 상용SW 4.0 P5 거부 시), 컨틴전시 소진 후 잔여.
\item \textbf{검토} — 재경팀장이 집행 한도·전결 규정·연차 자금 계획과의 정합을 3영업일 내 확인(RACI 20행 C).
\item \textbf{결재} — 정미란 스폰서 10영업일 내 결정(RACI 20행 A). 5억 초과 건은 프로그램 운영위 보고 병행, 프로그램 관리예비비 24억 총액 관리는 프로그램 PMO장.
\item \textbf{보고} — CCB 다음 회의에 보고, 윤태호 감사팀장 사본(2020 감사 지적 “집행 승인 근거” 대응).
\item \textbf{기준선 갱신} — 승인 금액만큼 계획 원가 상향(BL-0n), PV 곡선 해당 월 갱신, 집행 한도 동액 상향. 재기준선이 아니라 승인된 변경으로 기록.
\item \textbf{기록} — 집행 요청서·결재·CCB 보고를 형상 저장소에 보관, 종료 보고서 원가 절에 집행 이력 전량 첨부.
\end{enumerate}\fi
\end{fieldbox}

% ------------------------------------------------------------------ 8
\secbar[80mm]{8. 벤더 견적 검산}
\guide{주요 계약 금액을 두 방향에서 검산 — 규모 × 단가(FP × 단가)와 자원 × 기간(인력 × 개월 × 단가). 둘의 차이가 무엇(간접비·이윤·피크 인력)인지 적는다. 수행사의 목표 손익률과 대조해 변경 협상에서 수행사가 어디에 서 있는지 이해한다. 한 방향만 검산하거나 차이를 “이상하다”로 끝내지 말 것.}
{\footnotesize
\begin{tabularx}{\linewidth}{|L{24mm}|Y|C{16mm}|Y|}\hline
\hd{방향} & \hd{산식} & \hd{값 (억)} & \hd{차이·원인} \\ \hline
\ifwbblank
\rd{12mm}규모 × 단가 & \g{[FP × 단가]} & & \g{[계약 금액과 일치?]} \\ \hline
\rd{12mm}자원 × 기간 & \g{[인력 × 개월 × 단가 / 스프린트 n회 × 회당 비용]} & & \g{[차이 = 이윤 + 간접비 + 피크 증분]} \\ \hline
\rd{14mm}차이 분석 & \g{[차이 = 목표 손익률 × 계약 + 간접비·피크 …]} & & \g{[응찰가 → 계약가 인하가 어디에서 흡수되었나]} \\ \hline
\rd{10mm}\g{[단말 등 기타]} & & & \\ \hline
\rd{10mm}\g{[상용SW]} & & & \\ \hline
\else
규모 × 단가 & 12,400 FP × 550,000원 = 6,820,000,000원 & \textbf{68.20} & 계약 금액과 일치. L2 배분 3,100 / 2,900 / 2,300 / 1,500 / 2,600 × 0.55 = 17.05 / 15.95 / 12.65 / 8.25 / 14.30 \\ \hline
자원 × 기간 & 스프린트 36회 × 1.408억 (44명 × 0.5개월 × 640만 원, P-11a) & \textbf{50.69} & 평균 44명 기준 직접 인건비. 계약 68.20과의 차이 17.51 \\ \hline
차이 분석 & 17.51 = 목표 이윤 6.5\% × 68.2 ≈ 4.43 + 피크 61명 증분·간접비·위험 프리미엄 약 13.08 (68.2의 19\%) [추가 설정] & & 응찰가 198 → 계약 168(P-13)의 30억 인하가 이 19\% 안에서 흡수되었다는 뜻. 수행사의 변경대가 유인(R-05·R-07)은 여기서 나온다 \\ \hline
단말 & 477대 × 148만 원 = 7.06 + 설치 1.54 = 8.60 & 8.60 & 일치 \\ \hline
상용SW & 6.8 + 9.2 + 4.1 + 3.3 = 23.40 → 기준선 19.40 (2·3년차 유지보수 4.0 이관) & 19.40 & 벤더별 유지보수 요율 확인 필요(라이선스의 연 11~12\% 가정) \\ \hline
\fi
\end{tabularx}}

% ------------------------------------------------------------------ 9
\secbar[70mm]{9. 원가 변경 규칙}
\guide{원가 기준선을 바꾸는 트리거·절차·허용범위. 허용범위 격자(⑩ 항목 7)의 원가 축과 같은 숫자, \textbf{분모를 명시}(승인 예산의 n\%인가 계획 원가의 n\%인가). “컨틴전시 집행은 기준선 변경이 아니다(집행 한도 안의 이동), 관리예비비 집행과 범위 변경은 기준선 변경이다”를 명시한다. 컨틴전시를 쓸 때마다 기준선을 고치지 말 것.}
\begin{fieldbox}
\ifwbblank
\g{프로젝트 계층 원가 허용범위 (분모 명시): \hrulefill}\par\vspace{4mm}
\g{컨틴전시 집행 ≠ 기준선 변경 / 관리예비비·범위 변경 = 기준선 변경 — 문안: \hrulefill}\par\vspace{4mm}
\g{사내 전결·내부통제 지침과의 관계 (단일 건 기준 금액): \hrulefill}\par\vspace{4mm}
\g{초과 예상 판정 근거(EAC)와 보고 시한: \hrulefill}\par\vspace{4mm}
\g{재기준선 조건·승인·횟수: \hrulefill}\par
\else
\begin{itemize}
\item 프로젝트 계층 원가 허용범위 = \textbf{컨틴전시 10.2 전액 = 계획 원가의 +7.0\%} (Day1 초안 “+6\% = 10.08, 승인 예산 168 기준”에서 변경. 이유: 허용범위는 PM이 스폰서 보고 없이 쓸 수 있는 금액이어야 하고, 그것이 곧 컨틴전시다. 계획 원가 + 허용범위 = 집행 한도 156.0으로 층이 맞물린다 — ⑩ 항목 7·8).
\item 컨틴전시 집행은 기준선 변경이 아니다(집행 한도 안의 이동). 관리예비비 집행과 범위 변경(CCB 승인)은 기준선 변경이며 BL 번호가 올라간다.
\item 단일 건 5억 초과는 허용범위 내라도 CFO 병행 결재(헌장 §12), 5,000만 원 이상 발주는 내부통제 지침(2026-02-16).
\item 초과 \textbf{예상} 시점(EAC 산정 결과 집행 한도 초과 전망)에 5영업일 내 예외 보고. Day3 EVM의 EAC가 그 예상의 근거.
\item 재기준선은 CCB 승인 + 프로그램 이사회 확인, 최대 2회. 재기준선 시 PV 곡선 전체를 다시 발행하고 이전 곡선을 보관한다.
\end{itemize}\fi
\end{fieldbox}
\end{document}
